\documentclass[dvipdfmx]{jarticle}
\pagestyle{empty}
\usepackage{tikz}
\usetikzlibrary{shapes,calc}
\begin{document}
\begin{center}
  \begin{tikzpicture}
  \draw (0,0) rectangle (0.5,0.2);
  \draw (0.5,0) rectangle (1,0.2);
  \draw (1,0) rectangle (1.5,0.2);
  \draw (1.5,0) -- (2,0);
  \draw (1.5,0.2) -- (2,0.2);
  \draw (0,3.5) rectangle (0.5,3.7);
  \draw (0.5,3.5) rectangle (1,3.7);
  \draw (1,3.5) rectangle (1.5,3.7);
  \draw (1.5,3.5) -- (2,3.5);
  \draw (1.5,3.7) -- (2,3.7);
  \draw (0,1.75) rectangle +(2.5,1);
  \node at (1.3,2.25) {\parbox{2.5cm}{\scriptsize 有限制御部\\(遷移規則の集合)}};
  \draw (3.5,0.5) rectangle +(1,1.5);
  \draw (3.5,0.75) -- (4.5,0.75);
  \draw (3.5,1) -- (4.5,1);
  \draw (3.5,1.75) -- (4.5,1.75);
  \node at (4,1.5) {$\vdots$}; 
  \draw (3.5,2) -- (3.5,2.25); \draw (4.5,2) -- (4.5,2.25);
  \path[->] (4,2.5) edge (2.5,2.5)
            (4,2.5) edge (4,2.1);
  \node at (1,3.9) {\tiny 入力テープ};
  \node at (1,-0.2) {\tiny 出力テープ};
  \draw (1.25,2.75) -- (1.25,3);
  \draw (1.25,1.75) -- (1.25,1.25);
  \path[->] (1.25,1.25) edge (0.75,0.2);
  \path[->] (1.25,3) edge (0.25,3.5);
  \node[rotate=90] at (4.75,1.5) {\tiny プッシュダウンスタック};
  
  \end{tikzpicture}
\end{center}
\end{document}
